\section{Construction, the two-family selector, and finite chambers}

\subsection{Imported construction and normalized coordinates}
\label{sec:base}

Let $\A$ be an alphabet of cardinality $q=2m$, with a fixed-point-free
involution $a\mapsto\bar a$.  For a word $w=w_0\cdots w_{k-1}$, define
\[
  \RC(w)_j=\overline{w_{k-1-j}}.
\]
The directed vertex de Bruijn graph $B(q,k)$ has vertex set $\A^k$ and an arc
$u\to v$ when $u_1\cdots u_{k-1}=v_0\cdots v_{k-2}$.  A decycling set is a
set of vertices whose deletion leaves a directed acyclic graph.  An ordinary
minimum decycling set is abbreviated to MDS.

Rotation $R(w_0\cdots w_{k-1})=w_1\cdots w_{k-1}w_0$ partitions $\A^k$ into
pure cycling registers (PCRs).  Their number is the necklace count
\begin{equation}
  N_q(k)=\frac1k\sum_{d\mid k}\varphi(d)q^{k/d}.
  \label{eq:necklace-count}
\end{equation}
Every MDS has size $N_q(k)$ and contains exactly one distinct representative
from each PCR~\cite{Mykkeltveit1972,ChamparnaudHanselPerrin2004}.  The converse
is false: a one-per-PCR set can leave a cycle that traverses multiple PCRs.

We use the line-graph representation.  A length-$k$ word is the labelled arc
\[
  \tail(w)=w_0\cdots w_{k-2}\longrightarrow
  \head(w)=w_1\cdots w_{k-1}
\]
of $B(q,k-1)$.  For an integer function $f$ on lower states, set
\begin{equation}
  (\nabla f)(w)=f(\tail(w))-f(\head(w)).
  \label{eq:gradient}
\end{equation}

Let $k\ge4$ be even.  Choose distinct positive magnitudes
$\beta_1,\ldots,\beta_m$ and label each complement pair so that
\[
  s(a_i)=\beta_i,\qquad s(\bar a_i)=-\beta_i.
\]
Thus $s$ is real, injective, nonzero, and antisymmetric.  Put
$\zetaK=e^{2\pi i/k}$ and
\begin{align}
  P_s(w)&=\sum_{j=0}^{k-1}s(w_j)\zetaK^{j+1},                            \label{eq:embedding}\\
  W_s(w)&=\sum_{j=0}^{k-1}s(w_j),                                        \label{eq:weight-sum}\\
  \lambda_s(u)&=\sum_{j=0}^{k-2}s(u_j)
       \sin\!\left(\frac{2\pi(j+1)}k\right).                            \label{eq:lambda}
\end{align}
The weighted seed $M_0$ selects, on every nonzero embedded PCR, the negative
real vertex if present and otherwise the unique strict lower-to-upper
crossing.  On a zero-embedding PCR it uses the fixed deterministic
reverse-complement-equivariant tie rule of~\cite{Babanskyy2026Parity}.

For a lower word $u$, define
\begin{align}
  \psi_s(u)&=\sum_{j=0}^{k-2}s(u_j)(\zetaK^{j+1}-1),                     \label{eq:psi}\\
  r_s(u)&=\operatorname{Re}\psi_s(u),                                    \label{eq:midpoint}\\
  U_s&=\{u:\lambda_s(u)>0\},\qquad \vartheta(u)=\RC(u).                  \label{eq:positive-support}
\end{align}
Fix a lexicographic order on lower words and put
\begin{equation}
  \begin{aligned}
  a_s(u)=1\quad\Longleftrightarrow\quad&
  u\in U_s\ \text{and}\\
  &\bigl(r_s(u)<0\ \text{or}\ [r_s(u)=0\ \text{and}\ u<\vartheta(u)]\bigr).
  \end{aligned}
  \label{eq:as}
\end{equation}

We use the following theorem as an imported interface.

\begin{theorem}[Imported parity-selector interface]
\label{thm:imported-selector}
For every fixed-point-free complemented alphabet and every even $k\ge4$,
the seed $M_0$ is an ordinary MDS.  The function
\begin{equation}
  m_s(w)=m_0(w)+a_s(\tail(w))-a_s(\head(w))                               \label{eq:terminal-selector}
\end{equation}
is the indicator of a literal reverse-complement-invariant ordinary MDS.
The construction and proof require only real linearity, injectivity,
nonvanishing, and antisymmetry of $s$.
\end{theorem}

The seed statement is Theorem~4.3, the weighted selector is Theorem~6.3,
and its literal reverse-complement invariance and minimum-cardinality
consequences are Theorem~7.2 and Corollary~7.3 of
\cite{Babanskyy2026Parity}.  Their legal-descent and acyclicity proofs are not
repeated here.  We use only the formulas above and the fact that every
weight-dependent branch is a sign or zero test of an explicitly named real
linear form.

\subsubsection{Residual length convention}
Throughout, $L(M)$ counts edges in $B(q,k-1)-M$, viewing selected
length-$k$ words as arcs. Its line graph is $B(q,k)\setminus M$.
When $L(M)\ge1$, the longest upper residual path has $L(M)-1$ edges.
A longest word avoiding selected $k$-mers has $L(M)+k-1$ letters.

\subsubsection{Component notation}

For each letter write
\[
  s(x)=\sum_{i=1}^m \beta_i\varepsilon_i(x),
  \qquad \varepsilon_i(x)\in\{-1,0,1\},
\]
with exactly one nonzero component.  Define
\begin{align}
  X_i(w)&=\sum_{j=0}^{k-1}\varepsilon_i(w_j)\zetaK^{j+1},                 \label{eq:Xi}\\
  L_i(u)&=\sum_{j=0}^{k-2}\varepsilon_i(u_j)
    \frac{\sin(2\pi(j+1)/k)}{\sin(2\pi/k)}.                             \label{eq:Li}
\end{align}
Then
\begin{equation}
  P_s(w)=\sum_i\beta_iX_i(w),\qquad
  \lambda_s(u)=\sin(2\pi/k)\sum_i\beta_iL_i(u).                         \label{eq:component-decomp}
\end{equation}
If $K_k=\Q(\zetaK)$ and $K_k^+=\Q(\zetaK+\zetaK^{-1})$, then
$X_i(w)\in K_k$ and $L_i(u)\in K_k^+$.  The normalization by the common
positive sine in~\eqref{eq:Li} is essential; the sine itself need not lie in
the real cyclotomic subfield used for the coefficient test.

\subsection{Only two families are needed for the terminal}\label{ext:sec:two}
Write $V_m=\{\pm e_1,\ldots,\pm e_m\}$ for the signed one-hot
alphabet. For DNA take $T=e_1,A=-e_1,G=e_2,C=-e_2$ and
$\pi_\rho(h,l)=\rho h+l$. Structural vanishing means componentwise
vanishing. All path lengths below count edges of the lower residual
$D_s=B(q,k-1)-M_s$.
Let $w:u\to v$ be a lower-graph arc, $\alpha=2\pi/k$, and $b(w)$ the fixed zero-PCR seed membership rule, used only when $P_s(w)=0$.

\begin{theorem}[Two-family selector]\label{ext:thm:two}
For every admissible weight map and even $k\ge4$,
\[
 m_0(w)=
 \begin{cases}
 b(w),&\lambda_s(u)=\lambda_s(v)=0,\\
 1,&\lambda_s(u)\le0<\lambda_s(v),\\
 0,&\text{otherwise}.
 \end{cases}
\]
Consequently the signs of the two lower-word families $\lambda_s,r_s$, together with the frozen combinatorial tie rules, determine the seed, midpoint ideal, terminal, and labelled residual graph.
\end{theorem}
\begin{proof}
The suffix letter has exponent $k$ and contributes no imaginary part, so
\[
\lambda_s(u)=\Im P_s(w),\qquad
\lambda_s(v)=\Im(e^{-i\alpha}P_s(w)).
\]
Writing $P_s(w)=x+iy$, the second coordinate equals $y\cos\alpha-x\sin\alpha$. Since $\sin\alpha>0$, both lower coordinates vanish exactly when $x=y=0$. If $y=0$ and the embedding is nonzero, the next imaginary coordinate is positive exactly when $x<0$. This is the negative-real seed branch. Otherwise the seed is the strict lower-to-upper crossing, giving the displayed rule. The midpoint indicator remains
\[
a_s(u)=1\iff\lambda_s(u)>0\ \text{and}\
[r_s(u)<0\ \text{or}\ (r_s(u)=0\text{ and }u<\RC(u))].
\]
Substitute into $m_s(w)=m_0(w)+a_s(u)-a_s(v)$.
\end{proof}
This is a sufficient decision interface, not a claim that arbitrary raw seed-real signs can be recovered everywhere. The five-role certificate catalogue remains useful for classification and audit. The five-role trace is not an irredundant decision interface: the identities above supply the explicit aliases.

\begin{corollary}[Symmetry compression]\label{ext:cor:compression}
Let $Z^{\rm str}_\lambda$ count globally structural lower-lambda words. The terminal is determined by at most
\[
\frac{q^{k-1}-Z^{\rm str}_\lambda}{2}
\]
labelled real-form occurrences, in addition to the structural registry and fixed ties.
\end{corollary}
\begin{proof}
Let $C$ complement letters and $J$ reverse positions. The pairs $(\lambda,r)$ transform under $1,C,J,CJ$ with signs $(+,+),(-,-),(-,+),(+,-)$ respectively. A $J$-fixed lower word is palindromic and has componentwise zero lambda. Neither $C$ nor $CJ$ fixes an odd-length word over a fixed-point-free complemented alphabet. Thus the nonstructural-lambda lower words form free four-element orbits. Retain $\lambda$ and $r$ at one representative per orbit and restore all signs. A globally structural-lambda word always has $a_s=0$, so its midpoint sign is unnecessary.
\end{proof}
For DNA at $k=4,6,8$, the bound is respectively $24,440,7904$. These are occurrence counts before proportional-form deduplication, not minimum hyperplane counts.

\subsection{Complete decision covectors and terminal chambers}
\label{sec:covectors}

At fixed $\A$ and $k$, the selector has finitely many labelled branch
occurrences.  We retain the occurrence tags rather than deduplicating
coincident linear forms.  This makes branch completeness auditable.

\begin{definition}[Complete decision catalog]
\label{def:decision-catalog}
The catalog $\mathcal D_{\A,k}$ contains every weight-linear real form tested
by the deterministic construction, with its labelled word, rotation, and
branch role.  The roles are:
\[
  \role{seed-real},\quad \role{seed-imag},\quad
  \role{embedding-zero},\quad \role{lower-lambda},\quad
  \role{lower-midpoint}.
\]
The embedding-zero role records both real coordinates needed to decide
$P_s(w)=0$.  A form that is identically zero in weight space is retained as a
structural tie and is not treated as a hyperplane.  The \emph{complete
covector} $\sigma(s)$ is the sign in $\{-,0,+\}$ of every nonidentically-zero
catalog occurrence, together with the structural-tie registry.
\end{definition}

The catalog can be exponentially large in $k$; this section claims
finiteness and completeness, not a small irredundant arrangement.

\begin{theorem}[Terminal sign-covector factorization]
\label{thm:terminal-covector}
Fix a complemented alphabet, even $k\ge4$, the deterministic zero-PCR rule,
and the lexicographic midpoint tie rule.  If two admissible weight maps $s,t$
have the same complete decision covector, then they have identical:
\begin{enumerate}
  \item seed indicators and zero-PCR representatives;
  \item positive lower support and midpoint ideal;
  \item terminal indicators $m_s=m_t$ word by word;
  \item labelled residual digraphs; and
  \item deterministic statistics computed from those labelled digraphs.
\end{enumerate}
Consequently the complete terminal payload is constant on every relatively
open face of the realizable arrangement defined by $\mathcal D_{\A,k}$ and is
invariant under positive scaling of all magnitudes.
\end{theorem}

\begin{proof}
The signs of the tagged seed-real and seed-imag occurrences determine the
negative-ray and strict-crossing branches on every nonzero PCR.  The paired
coordinates tagged embedding-zero determine exactly which PCRs use the
fixed, weight-independent zero rule.  Hence the seed indicators agree.

The lower-lambda signs determine the same set $U_s=U_t$.  On that set, the
lower-midpoint signs, followed only when zero by the fixed lexicographic rule,
determine the same function $a_s=a_t$.  Substitution in the endpoint formula
\eqref{eq:terminal-selector} gives $m_s=m_t$ on every labelled word.  Deleting
the same labelled vertices gives the same residual digraph, so any
deterministic statistic of it is identical.

The sign covectors are constant on relatively open faces of a finite real
hyperplane arrangement~\cite{BjornerEtAl1999}.  Every catalog form is
homogeneous of degree one in the magnitudes, so multiplication by a positive
scalar preserves all signs and all tie decisions.
\end{proof}

The five-role trace is sufficient but not irredundant.
\Cref{ext:thm:two} reduces it to two lower-word families with fixed ties.
Equal aggregate zero counts alone still do not determine the terminal.

\begin{remark}[Resonance strata versus terminal strata]
\label{rem:resonance-terminal-strata}
The resonance stratification records only which selector equalities vanish.
The terminal decision arrangement records the full sign of every labelled
catalog occurrence.  It therefore refines the resonance stratification:
equal zero loci need not determine equal seed choices, midpoint ideals, or
terminals.  The bounded atlas of \Cref{prop:finite-atlas} exhibits strict
refinement in its declared rows; no all-order irredundancy claim is inferred
from that finite witness.
\end{remark}

\begin{corollary}[Exact one-parameter DNA diagram]
\label{cor:dna-phase}
For DNA weights
\begin{equation}
  s(A)=-\rho,\quad s(C)=-1,\quad s(G)=1,\quad s(T)=\rho,
  \qquad \rho>1,                                                        \label{eq:dna-weights}
\end{equation}
every catalog form is $a\rho+b$ with real algebraic coefficients $a,b$.
After structural forms are removed, the exact phase diagram is the ordered
finite set of realized critical ratios $-b/a>1$, the intervening open
intervals, and the critical points evaluated with the frozen tie rules.  The
complete terminal payload is constant on every such cell.
\end{corollary}

\begin{proof}
Every selector form is real-linear in the weight table.  Substitution of
\eqref{eq:dna-weights} leaves an affine function of $\rho$.  A nonzero affine
form can change sign only at its unique root.  Exact algebraic comparison
orders the finitely many admissible roots.  The conclusion follows from
\Cref{thm:terminal-covector}; point cells are kept separate because their
zero branches need not agree with either adjacent interval.
\end{proof}

This corollary proves a complete diagram only after the entire catalog has
been generated.  Omitting a branch family can merge cells whose terminals
are different.
